\chapter{Cadre général du projet}
\label{chap:cadre-general}

\section*{Introduction}
\addcontentsline{toc}{section}{Introduction}

Ce premier chapitre situe le stage dans son contexte : l'entreprise d'accueil, la problématique à l'origine du projet, les objectifs fixés et la démarche méthodologique retenue pour y répondre.

\section{Présentation de l'organisme d'accueil}

\textbf{Prodexo} est l'entreprise au sein de laquelle ce stage de fin d'études a été réalisé. \emph{[Section à compléter avec la présentation officielle de l'entreprise : secteur d'activité, taille, positionnement, principaux clients ou domaines d'intervention, organigramme du service d'accueil.]}

Le stage s'est déroulé au sein de l'équipe en charge du développement produit, sur un projet interne de plateforme d'assistance IA destinée aux marchands utilisant des solutions e-commerce, en particulier \textbf{PrestaShop}.

\section{Contexte du projet}

À l'arrivée sur le projet, une première version de la plateforme existait déjà : un service backend en Python/FastAPI capable de répondre à des messages de chat en s'appuyant sur un modèle de langage, une intégration partielle avec PrestaShop pour la synchronisation du catalogue produits, ainsi qu'une ébauche de backoffice en Laravel/Filament destinée aux administrateurs de boutique.

Cette première version souffrait cependant de deux limites majeures, identifiées dès la phase de prise en main :

\begin{itemize}
    \item \textbf{Un comportement du chat figé dans le code.} Toute personnalisation du comportement de l'assistant (réponses automatiques à certains mots-clés, génération de coupons de fidélisation, politiques de remise) nécessitait une modification du code source et un redéploiement, ce qui rendait la plateforme peu exploitable par un administrateur non technique.
    \item \textbf{Des données d'administration partiellement fabriquées.} Plusieurs écrans du backoffice -- tableau de bord, agent IA d'administration, recommandations -- affichaient des valeurs d'exemple codées en dur (chiffre d'affaires fictif, noms de produits génériques) plutôt que des données réellement calculées à partir de l'activité de la plateforme. Ce constat, documenté au fil de l'audit du code existant, a orienté une part importante du projet vers un travail de \emph{fiabilisation} plutôt que de seule addition de fonctionnalités.
\end{itemize}

\section{Problématique}

La problématique du stage peut se formuler ainsi : \emph{comment permettre à un administrateur de boutique de configurer et de superviser un assistant IA e-commerce en toute confiance, sans dépendre d'une intervention technique pour chaque ajustement, et sans que les informations qui lui sont présentées soient trompeuses ?}

Cette problématique se décline en plusieurs sous-questions techniques traitées tout au long du projet :

\begin{itemize}
    \item Comment permettre la personnalisation du comportement du chat (réponses automatiques, remises, instructions contextuelles) sans redéploiement, tout en gardant un système simple à auditer ?
    \item Comment donner à un agent IA la capacité d'exécuter des actions d'administration (génération de coupons en masse, mise à jour de prix) sans lui accorder une exécution libre et potentiellement dangereuse ?
    \item Comment calculer des indicateurs de pilotage (produits les plus demandés, demande non satisfaite, taux de conversion des coupons) à partir des seules données réellement disponibles, en l'absence -- notable -- de table de commandes dans le système ?
    \item Comment garantir, via des tests automatisés reproductibles, que ces garanties (isolation multi-tenant, sécurité de l'agent d'administration, exactitude des métriques) restent vraies dans la durée ?
\end{itemize}

\section{Objectifs du projet}

Sur la base de cette problématique, les objectifs suivants ont été fixés en début de stage :

\begin{enumerate}
    \item Concevoir et implémenter un \textbf{moteur de règles métier} permettant à un administrateur de définir, via une interface graphique, des conditions (intention détectée, mots-clés) déclenchant des actions (réponse automatique, instruction injectée dans le contexte du modèle, génération de coupon), interprétées par le chat sans modification de code.
    \item Remplacer les données fabriquées de l'agent IA d'administration par des calculs réels, et sécuriser son exécution par un mécanisme d'actions prédéfinies avec confirmation.
    \item Concevoir un service d'\textbf{analyse de la demande client} (\emph{insights}) fondé exclusivement sur les données réellement collectées par la plateforme.
    \item Faire évoluer le backoffice Laravel/Filament pour exposer ces nouvelles capacités : gestion des règles, interface de pilotage de l'agent IA, tableaux de bord à données réelles, suivi des coûts d'utilisation du modèle de langage.
    \item Couvrir l'ensemble de ces évolutions par une suite de tests automatisés et documenter l'architecture finale du système.
\end{enumerate}

\section{Méthodologie de gestion de projet}

\subsection{Choix de la méthodologie}

Le projet a été mené selon une démarche agile inspirée du cadre \textbf{Scrum}~\cite{scrumguide}. Ce choix se justifie par la nature du projet : un périmètre fonctionnel amené à être précisé au fil de l'avancement (l'audit de l'existant a lui-même fait émerger une partie du périmètre réel du stage), la nécessité de livrer des incréments testables et démontrables régulièrement, et un stage solo qui bénéficie d'un cadre stagiaire/encadrant proche d'un binôme Product Owner / équipe de développement.

Le déroulement du stage a ainsi suivi les grands principes de Scrum, adaptés à un contexte de stage individuel :

\begin{itemize}
    \item un \textbf{backlog produit} recensant les besoins fonctionnels, priorisé conjointement avec l'encadrant professionnel ;
    \item un découpage du développement en \textbf{sprints} de deux semaines, chacun clos par une démonstration des fonctionnalités livrées ;
    \item une revue de sprint informelle avec l'encadrant en fin de période, permettant d'ajuster le backlog du sprint suivant ;
    \item une exigence de \emph{définition of done} stricte pour chaque fonctionnalité : tests automatisés passants (unitaires et, lorsque pertinent, d'intégration sur base de données réelle), vérification manuelle de l'interface concernée, et documentation à jour.
\end{itemize}

\subsection{Découpage temporel du stage}

Le stage de six mois s'est organisé en deux grandes périodes, détaillées au \cref{chap:realisation} :

\begin{itemize}
    \item une \textbf{phase de cadrage et de conception} (février à avril 2026), consacrée à la prise en main de l'existant, à la formalisation du cahier des charges, à la montée en compétence sur la pile technique (FastAPI, SQLAlchemy asynchrone, Filament) et à la conception de l'architecture cible ;
    \item une \textbf{phase de réalisation} (mai à août 2026), structurée en six sprints de développement itératifs, décrits en détail au chapitre 4.
\end{itemize}

\begin{table}[h]
\centering
\caption{Planification générale du stage}
\label{tab:planification}
\begin{tabular}{@{}p{4cm}p{3cm}p{7cm}@{}}
\toprule
\textbf{Période} & \textbf{Durée} & \textbf{Contenu} \\
\midrule
Cadrage \& conception & Fév. -- Avr. 2026 & Analyse de l'existant, spécification des besoins, conception de l'architecture \\
Sprint 1 & Mai (S1--S2) & Fondations \& infrastructure \\
Sprint 2 & Mai (S3--S4) & Moteur de chat IA (RAG, persistance) \\
Sprint 3 & Juin (S1--S2) & Sécurité, multi-tenant \& supervision \\
Sprint 4 & Juin (S3--S4) & Backoffice Laravel/Filament -- socle \\
Sprint 5 & Juillet (S1--S2) & Moteur de règles \& agent IA d'administration \\
Sprint 6 & Juillet -- Août & Finalisation backoffice, analytique réelle, documentation \\
\bottomrule
\end{tabular}
\end{table}

\section{Environnement de travail}

Le projet a été développé dans un environnement conteneurisé (\textbf{Docker} / Docker Compose), avec un dépôt \textbf{Git} unique hébergeant les deux composants applicatifs (\emph{ai-core} et \emph{backoffice}) ainsi que l'infrastructure de déploiement. L'intégration continue s'appuie sur \textbf{GitHub Actions}, avec analyse de qualité de code (\emph{linting}, formatage) exécutée avant chaque fusion. Le suivi de version a suivi une convention de \emph{commits} typés (\texttt{feat}, \texttt{fix}, \texttt{chore}) facilitant la traçabilité du travail réalisé sprint par sprint.

\section*{Conclusion}
\addcontentsline{toc}{section}{Conclusion}

Ce chapitre a présenté le contexte, la problématique et les objectifs du stage, ainsi que la démarche méthodologique adoptée pour y répondre. Le chapitre suivant détaille l'analyse et la spécification des besoins qui ont découlé de cette problématique.
